
\begin{figure}[H]
\centering
\begin{tikzpicture}[
  scale=0.95, transform shape,
  >=Latex,
  every node/.style={font=\footnotesize},
  mainwidth/.store in=\mainwidth, mainwidth=12.3cm,
  box/.style   ={rectangle, rounded corners=2mm, draw, very thick, fill=gray!5,
                 align=center, inner sep=2mm, text width=\mainwidth}, % zentriert
  io/.style    ={box, fill=orange!15},                                 % erbt zentriert
  proc/.style  ={box, fill=blue!8},                                    % erbt zentriert
  decision/.style={diamond, draw, very thick, aspect=2, fill=green!10,
                   inner sep=1.0mm, align=center,
                   minimum width=4.0cm, minimum height=1.0cm},
  arr/.style   ={->, very thick},
  arrthin/.style={->, line width=0.8pt},
  loopframe/.style={draw=black!60, rounded corners=3mm, very thick, dashed}
]

% --- Spaltenlayout mit definierter vertikaler Distanz ---
\matrix[column sep=0mm, row sep=7mm] (m) {
  \node (start) [box] {Start Client (\texttt{python client.py})}; \\
  \node (sess)  [proc]{Session \& Logging}; \\
  \node (env)   [box] {Konfiguration: Preprocessing-Parameter, CAE-Architektur, Strategy}; \\
  \node (load)  [io]  {CSV aus MinIO laden}; \\
  \node (prep)  [proc]{Vorverarbeitung pro Operation}; \\
  \node (win)   [proc]{Sliding Windows \& Splits (Train/Val/Test)}; \\
  \node (model) [proc]{Modellbau (CAE)}; \\
  \node (strat) [decision]{\textbf{STRATEGY}}; \\
  % ---- Federated-Learning-Loop ----
  \node (local)    [proc]{Lokales Training (Client)}; \\
  \node (exchange) [io]  {Modellaustausch mit Server}; \\
  \node (eval)     [proc]{Evaluation}; \\
  \node (save)     [io]  {Artefakte \& Uploads nach MinIO (Metriken, Logs; Modelle)}; \\
  \node (more)     [decision]{Weitere Runde?}; \\
  % ---- Ende ----
  \node (end)      [box] {Ende (Client beendet)}; \\
};

% Rahmen um den FL-Loop
\node[loopframe, fit=(local)(exchange)(eval)(save)(more),
      inner xsep=6mm, inner ysep=6mm,
      label={[anchor=south east, xshift=-31mm, yshift=1mm]\footnotesize Federated-Learning-Loop}] (loopbox) {};

% Fluss-Pfeile
\draw[arr] (start) -- (sess);
\draw[arr] (sess)  -- (env);
\draw[arr] (env)   -- (load);
\draw[arr] (load)  -- (prep);
\draw[arr] (prep)  -- (win);
\draw[arr] (win)   -- (model);
\draw[arr] (model) -- (strat);
\draw[arr] (strat) -- (local);
\draw[arr] (local) -- (exchange);
\draw[arr] (exchange) -- (eval);
\draw[arr] (eval) -- (save);
\draw[arr] (save) -- (more);

% Loop zurück (JA) bzw. Ende (NEIN)
\coordinate (retA) at ([xshift=-56mm]more.west);
\coordinate (retB) at ([xshift=-10mm]local.west);
\draw[arr] (more.west) -- node[pos=0.08, above]{\footnotesize Ja} (retA)
                        -- (retB) -- (local.west);
\draw[arr] (more.south) -- node[right]{\footnotesize Nein} (end.north);

% Callouts (dünne Pfeile) Richtung Server
\coordinate (rightlane) at ([xshift=10mm]loopbox.east);
\draw[arrthin] (local.east) -- (rightlane |- local.east)
  node[pos=1, right]{\scriptsize Fit-Metriken $\rightarrow$ Server};
\draw[arrthin] (eval.east) -- (rightlane |- eval.east)
  node[pos=1, right]{\scriptsize Eval-Metriken $\rightarrow$ Server};

\end{tikzpicture}
\caption{Client-Workflow mit Vorverarbeitung, Fensterung, Modellbau und Federated-Learning-Loop (Austausch \& Metrik-Übermittlung via Flower).}
\label{fig:client-workflow}
\end{figure} \pagebreak